\section{Introduction}
\label{sec:intro}

\subsection{The problem}
A rotating savings and credit association (ROSCA), called a \emph{committee} or \emph{kameti} in
Pakistan, is a group of people who each pay the same amount every round. In every round one member
receives the whole pot, until everyone has had a turn. Members paid early receive an interest free
loan; members paid late are saving by force. Committees are widely used by low income groups in
Pakistan \cite{mehmood2019}.

Everything rests on one person, the organizer, who
\begin{itemize}
  \item \textbf{holds the money} between collection and payout, and can delay, divert or disappear with it;
  \item \textbf{keeps the record} of who paid, usually in a notebook or a chat that members cannot check;
  \item \textbf{decides the payout order}, which is valuable (an early slot is a loan) and open to favouritism;
  \item \textbf{absorbs or hides defaults}: when a member who already took the pot stops paying, the
        members paid late simply lose money.
\end{itemize}

\subsection{What we built}
\emph{Trustless ROSCA} is an Ethereum decentralised application (DApp) that removes those four powers
from the organizer. Table~\ref{tab:powers} summarises how.

\begin{table}[H]
\centering\small
\caption{The organizer's powers and what replaces them}
\label{tab:powers}
\begin{tabularx}{\textwidth}{@{}P{3.3cm}L@{}}
\toprule
Organizer power today & In Trustless ROSCA \\
\midrule
Holds the pot & Contributions sit in the \code{Circle} contract and deposits in \code{CollateralVault}. No function lets the organizer, the admin or the developers move anyone else's money. \\
Keeps the record & Every join, payment, default and payout is a public event that every member can read. \\
Chooses the order & \code{PayoutOrder} runs a commit and reveal draw; nobody, including the last member to reveal, can pick a slot. \\
Absorbs defaults & Each member's deposit always covers 100\% of what they still owe after their payout, so a member who takes the pot and disappears is paid for out of their own deposit. \\
\bottomrule
\end{tabularx}
\end{table}

{\sloppy The system has seven contracts (\code{BaseCircleLifecycle}, \code{Circle},
\code{CircleFactory}, \code{CollateralVault}, \code{PayoutOrder}, \code{ReputationRegistry} and
\code{MockStablecoin}), four off chain services (IPFS terms store, event indexer with an HTTP API, in app notifier, keeper) and a
React frontend. A circle goes through eight user workflows (W1 to W8) across a five stage lifecycle,
with four roles: Organizer, Member, platform Admin (pause only) and anyone (keeper).\par}

\subsection{Why a blockchain, and not a bank or a payments app}
The obvious question is why not JazzCash, Easypaisa, a bank or a committee app.
\begin{itemize}
  \item \textbf{An app moves the trust problem; it does not remove it.} A company that holds the pot can
        freeze, lose or misuse it, can be hacked, and can change its rules. Our contract holds the pot,
        its rules cannot change after creation, and nobody can withdraw someone else's money.
  \item \textbf{Shared state between people who do not fully trust each other.} Members need one record of
        who paid that nobody can edit afterwards. A public chain gives every member the same
        verifiable ledger without a trusted record keeper.
  \item \textbf{Rules that run on deadlines.} Settlement, default handling and slashing happen in code
        when a deadline passes, triggered by anyone. No organizer needs to be honest or awake.
  \item \textbf{Verifiable fairness.} The payout order is derived from secrets every member contributed.
\end{itemize}
Blockchain is used only where it adds trust. Names, the terms document, reminders and history stay
off chain (Section~\ref{sec:datamap}).

\subsection{Objectives and scope}
The project brief \cite{brief} asks for a justified use case, a multi contract architecture with inheritance,
on chain and off chain integration, gas measurement with optimizations, a security analysis with
attack tests, a working frontend, tests and a report. Our objectives were therefore:
\begin{enumerate}
  \item hold every token in contracts and make defaults cost the defaulter, not the group;
  \item draw a payout order that no participant can choose;
  \item keep personal and bulky data off chain without making correctness depend on it;
  \item measure gas on the baseline and cut it with at least two optimizations;
  \item attack our own system, fix what the attacks expose and document what remains.
\end{enumerate}
Prototype choices: a mock stablecoin with 6 decimals (so values do not swing like ETH) on a local
Hardhat network. The guarantor role, Chainlink VRF, email or Telegram alerts and member matching
were cut to fit 14 days and are future work (Section~\ref{sec:limits}).

\subsection{How the work was organised}
The code moved as a relay over 14 days, one member at a time, each stage ending with a git tag and a
handoff note: Member A built the seven contracts and their unit tests (tag \code{handoff\_1});
Member B built the off chain services and the frontend on top of them (tag \code{handoff\_2});
Member C measured and optimized gas, attacked the system, fixed what the attacks exposed, wrote the
threat model and assembled this report (tag \code{final}). Section~\ref{sec:contrib} gives the
contribution matrix.

\subsection{Structure of this report}
Section~\ref{sec:related} reviews related work. Section~\ref{sec:design} presents the design:
roles, workflows, lifecycle, collateral rule, architecture, contracts and the data map.
Section~\ref{sec:impl} describes the implementation. Section~\ref{sec:gas} reports gas measurement
and the four optimizations, and Section~\ref{sec:security} the threat model, attack tests and the
three vulnerabilities found and fixed. Section~\ref{sec:testing} gives test results,
Section~\ref{sec:limits} limitations and future work, Section~\ref{sec:contrib} the work
distribution, and Section~\ref{sec:summary} a summary.
